\documentclass[12pt]{article}
\usepackage{amsmath, amssymb}
\usepackage[margin=1in]{geometry}
\setlength{\parskip}{6pt}
\title{QUBO Formulation: Molecular Conformation Selection Problem}
\date{}
\begin{document}
\maketitle

\subsection*{Molecular Conformation Selection Problem}

\paragraph{Problem Statement.}
Given a molecule decomposed into $N$ rotatable torsions, where torsion $i$ admits $M_i$ discrete conformational states (rotamers), the task is to select exactly one rotamer for each torsion such that the total molecular energy is minimized. The total energy comprises the sum of local energies associated with each selected rotamer and the sum of pairwise interaction energies between rotamers on distinct torsions.

\paragraph{Decision Variables.}
For each torsion $i \in \{0, \dots, N-1\}$ and each candidate rotamer $j \in \{0, \dots, M_i-1\}$, we introduce a binary decision variable $x_{i,j} \in \{0,1\}$. The variable $x_{i,j}$ equals $1$ if and only if rotamer $j$ is selected for torsion $i$, and $0$ otherwise.

\paragraph{QUBO Derivation.}
Step 1 --- Write the original objective in general notation. The minimization problem is:
\begin{align}
\min_{\mathbf{x}} \quad & f(\mathbf{x}) = \sum_{i=0}^{N-1} \sum_{j=0}^{M_i-1} E_{i,j} x_{i,j} + \sum_{i=0}^{N-1} \sum_{k=i+1}^{N-1} \sum_{j=0}^{M_i-1} \sum_{l=0}^{M_k-1} W_{i,j,k,l} x_{i,j} x_{k,l},
\end{align}
where $E_{i,j}$ denotes the local energy of rotamer $j$ on torsion $i$, and $W_{i,j,k,l}$ denotes the pairwise interaction weight between rotamer $j$ on torsion $i$ and rotamer $l$ on torsion $k$.

Step 2 --- Write each constraint in its original algebraic form. For every torsion $i$, exactly one rotamer must be chosen:
\begin{align}
\sum_{j=0}^{M_i-1} x_{i,j} = 1, \quad \forall i \in \{0, \dots, N-1\}.
\end{align}

Step 3 --- Convert each constraint to a penalty term. For each torsion $i$, we introduce a penalty term with weight $\lambda > 0$:
\begin{align}
P_i(\mathbf{x}) = \lambda \left( \sum_{j=0}^{M_i-1} x_{i,j} - 1 \right)^2.
\end{align}
Expanding the square yields:
\begin{align}
P_i(\mathbf{x}) = \lambda \left( \sum_{j=0}^{M_i-1} x_{i,j}^2 + \sum_{j=0}^{M_i-1} \sum_{\substack{l=0 \\ l \neq j}}^{M_i-1} x_{i,j} x_{i,l} - 2 \sum_{j=0}^{M_i-1} x_{i,j} + 1 \right).
\end{align}
Applying the idempotent property $x_{i,j}^2 = x_{i,j}$ for binary variables simplifies the expression to:
\begin{align}
P_i(\mathbf{x}) = \lambda \left( \sum_{j \neq l} x_{i,j} x_{i,l} - \sum_{j=0}^{M_i-1} x_{i,j} + 1 \right).
\end{align}
The total penalty contribution to the Hamiltonian is $\sum_{i=0}^{N-1} P_i(\mathbf{x})$.

Step 4 --- Combine objective and all penalty terms into the single QUBO Hamiltonian $H(\mathbf{x})$:
\begin{align}
H(\mathbf{x}) = \sum_{i=0}^{N-1} \sum_{j=0}^{M_i-1} E_{i,j} x_{i,j} + \sum_{i=0}^{N-1} \sum_{k=i+1}^{N-1} \sum_{j=0}^{M_i-1} \sum_{l=0}^{M_k-1} W_{i,j,k,l} x_{i,j} x_{k,l} + \sum_{i=0}^{N-1} \lambda \left( \sum_{j \neq l} x_{i,j} x_{i,l} - \sum_{j=0}^{M_i-1} x_{i,j} + 1 \right).
\end{align}

Step 5 --- Read off the Q matrix entries and offset $c$ from $H$ in general closed form. Grouping terms by variable pairs $(i,j)$ and $(k,l)$ gives:
\begin{align}
Q_{(i,j),(i,j)} &= E_{i,j} - \lambda, \quad \forall i \in \{0,\dots,N-1\}, \, j \in \{0,\dots,M_i-1\}, \\
Q_{(i,j),(i,l)} &= \lambda, \quad \forall i \in \{0,\dots,N-1\}, \, j,l \in \{0,\dots,M_i-1\}, \, j \neq l, \\
Q_{(i,j),(k,l)} &= W_{i,j,k,l}, \quad \forall i,k \in \{0,\dots,N-1\}, \, i < k, \, j \in \{0,\dots,M_i-1\}, \, l \in \{0,\dots,M_k-1\}, \\
Q_{(i,j),(k,l)} &= 0, \quad \text{otherwise}, \\
c &= \sum_{i=0}^{N-1} \lambda.
\end{align}

\paragraph{Penalty Weight Justification.}
The penalty weight $\lambda$ must be chosen sufficiently large to enforce feasibility. Specifically, $\lambda > \max_{i,j} |E_{i,j}|$ guarantees that the energy penalty for violating any single constraint strictly exceeds the maximum possible reduction in the objective function achievable by any feasible assignment. This ensures that infeasible configurations are energetically unfavorable compared to all feasible ones.

\paragraph{Correctness Argument.}
Minimizing $H(\mathbf{x})$ over $\{0,1\}^n$ recovers the optimal conformation because (i) for any feasible assignment, all constraint terms vanish, reducing $H(\mathbf{x})$ to the original objective function; (ii) for any infeasible assignment, at least one constraint is violated, contributing at least $\lambda$ to the total energy. Since $\lambda$ dominates the maximum possible objective gain, any infeasible solution has strictly higher energy than the optimal feasible solution. Consequently, the global minimum of $H(\mathbf{x})$ corresponds exactly to the feasible assignment that minimizes the original molecular energy.

\paragraph{Illustrative Example.}
Consider a molecule with $N=2$ torsions, where $M_0 = 2$ and $M_1 = 2$. The decision variables are $x_{0,0}, x_{0,1}, x_{1,0}, x_{1,1}$. The objective function becomes:
\begin{align}
f(\mathbf{x}) = E_{0,0}x_{0,0} + E_{0,1}x_{0,1} + E_{1,0}x_{1,0} + E_{1,1}x_{1,1} + W_{0,0,1,0}x_{0,0}x_{1,0} + W_{0,0,1,1}x_{0,0}x_{1,1} + W_{0,1,1,0}x_{0,1}x_{1,0} + W_{0,1,1,1}x_{0,1}x_{1,1}.
\end{align}
The penalty terms for the two torsions are:
\begin{align}
\lambda \left( x_{0,0} + x_{0,1} - 1 \right)^2 + \lambda \left( x_{1,0} + x_{1,1} - 1 \right)^2.
\end{align}
Expanding and applying $x_{i,j}^2 = x_{i,j}$ yields:
\begin{align}
\lambda \left( x_{0,0}x_{0,1} - x_{0,0} - x_{0,1} + 1 \right) + \lambda \left( x_{1,0}x_{1,1} - x_{1,0} - x_{1,1} + 1 \right).
\end{align}
The complete QUBO Hamiltonian is:
\begin{align}
H(\mathbf{x}) = (E_{0,0}-\lambda)x_{0,0} + (E_{0,1}-\lambda)x_{0,1} + (E_{1,0}-\lambda)x_{1,0} + (E_{1,1}-\lambda)x_{1,1} + 2\lambda x_{0,0}x_{0,1} + 2\lambda x_{1,0}x_{1,1} + W_{0,0,1,0}x_{0,0}x_{1,0} + W_{0,0,1,1}x_{0,0}x_{1,1} + W_{0,1,1,0}x_{0,1}x_{1,0} + W_{0,1,1,1}x_{0,1}x_{1,1} + 2\lambda.
\end{align}
The constant offset is $c = 2\lambda$. For any concrete instance parameters $\{E_{i,j}, W_{i,j,k,l}\}$, the optimal rotamer selection $\mathbf{x}^*$ is the binary vector minimizing $H(\mathbf{x})$, and the minimum total energy is $H(\mathbf{x}^*) - c$.

\end{document}